% !TeX program = xelatex
\documentclass[11pt,a4paper]{article}
\usepackage[margin=23mm]{geometry}
\usepackage{fontspec}
\setmainfont{Latin Modern Roman}
\usepackage{amsmath,amssymb,mathtools}
\usepackage{xcolor,enumitem,booktabs,longtable,array,xurl,hyperref}
\definecolor{muted}{HTML}{53636C}
\hypersetup{colorlinks=true,urlcolor=blue,linkcolor=blue}
\setlength{\parindent}{0pt}
\setlength{\parskip}{5pt}
\newcommand{\R}{\mathbb R}
\newcommand{\E}{\mathbb E}
\newcommand{\N}{\mathbb N}
\newcommand{\ind}{\mathbf 1}
\newcommand{\Var}{\operatorname{Var}}
\newcommand{\Cov}{\operatorname{Cov}}
\newcommand{\supp}{\operatorname{supp}}
\DeclareMathOperator*{\argmin}{arg\,min}
\DeclareMathOperator*{\argmax}{arg\,max}
\newcommand{\entrymeta}[3]{{\sffamily\small\color{muted}\textbf{\texttt{#1}}\quad|\quad #2\par #3\par}}
\newcommand{\Problemref}[1]{\texttt{#1}}
\newcommand{\JELref}[2]{\texttt{#2} (JEL \texttt{#1})}
\begin{document}
\section*{Parameterized computation of a reachable Pareto-optimal allocation}
\textbf{Problem ID:} \texttt{C78-369}\quad \textbf{JEL:} \texttt{C78}\par
% BEGIN ECONOMICS STATEMENT
\label{problem:OP-0552}
\label{beyond:MM-11}
\entrymeta{MM-11}{Matching markets and allocation}{Research agenda formalization}
\subsection*{Definitions and assumptions}
There are $n$ applicants $A$ and $n$ houses $H$, with $A$ and $H$ disjoint, and a simple undirected social network $G=(A,F)$. Each applicant $a$ has an acceptable set $B_a\subseteq H$ and a strict linear order $\succ_a$ on $B_a$; for $x,y\in B_a$, write $x\succeq_a y$ if $x=y$ or $x\succ_a y$. Allocations are bijections $M:A\to H$ with $M(a)\in B_a$ for all $a$; the endowment $M_0$ is such an allocation. A legal swap at $M$ exchanges the houses of distinct applicants $a,b$ precisely when $\{a,b\}\in F$, $M(b)\in B_a$, $M(a)\in B_b$, $M(b)\succ_a M(a)$ and $M(a)\succ_b M(b)$, leaving every other assignment unchanged. Let $\mathcal R_G(M_0)$ be all allocations obtainable through finite legal-swap sequences, including $M_0$. A reachable allocation $M$ is Pareto optimal relative to this set if
\[
 \forall N\in\mathcal R_G(M_0),\quad
 \bigl[\forall a\in A,\ N(a)\succeq_a M(a)\bigr]
 \Longrightarrow \bigl[\forall a\in A,\ N(a)=M(a)\bigr].
\]
Strict preferences make the equality formulation equivalent to the absence of a strict Pareto improvement.
\subsection*{Formal question}
Let $I$ consist of the explicitly encoded sets, graph, ordered acceptable lists and endowment, and let $|I|$ be their total bit length. Determine computable parameters $\kappa(I)$ for which one can output such an $M$, together with a legal sequence reaching it, in time $f(\kappa(I))|I|^c$, where $f$ is computable and $c$ is a constant independent of $I$ and the parameter value. Establish parameterized hardness for the corresponding search task or suitably specified decision formulations where appropriate; each such claim must specify its reduction convention and, for a decision formulation, its predicate and connection to the search task. Bare existence of a reachable Pareto-optimal allocation is not a nontrivial decision formulation. Candidate parameters include the treewidth of $G$, maximum list length $\ell=\max_{a\in A}|B_a|$, and the number of applicants permitted to exchange their houses. For the last candidate, it remains unspecified whether the applicants form a supplied eligible set, are selected subject to a participation budget, or are counted from an output or witness, and whether any restriction changes the Pareto comparison set. No such restriction is imposed on $\mathcal R_G(M_0)$ above; a formulation using this candidate must settle these choices explicitly.
\subsection*{Scope and status}
The 2019 survey proposes parameterized study of finding reachable Pareto-optimal allocations and reports hardness even on trees in Theorem~6. However, the cited 2017 paper proves general-network hardness in Proposition~4 and leaves the tree case open in Section~6. The treewidth-only FPT obstruction would follow from the survey's reported tree hardness, unless $\mathrm P=\mathrm{NP}$; the discrepancy between the cited texts remains unresolved here. Greedily continuing improving swaps until none is possible only produces a local terminal allocation; another branch from $M_0$ may produce a Pareto-superior reachable allocation. The global reachable-set definition is essential. Every finite reachable set has a Pareto-optimal member, so the obstacle concerns computation, not existence. The specific parameter combinations above are formalized research directions, with source-dated status.
\subsection*{References for the source of the problem}
Katar\'ina Cechl\'arov\'a, \'Agnes Cseh, and David F. Manlove (2019), \emph{Selected open problems in Matching Under Preferences}, Section~3.1.3, the reachable-set Pareto definition, Theorem~6 and concluding paragraph: \url{https://hpi.de/friedrich/docs/publications/2019/BEATCS.pdf}. Laurent Gourv\`es, Julien Lesca, and Ana\"elle Wilczynski (2017), \emph{Object Allocation via Swaps along a Social Network}, IJCAI 2017: \url{https://doi.org/10.24963/ijcai.2017/31}.

\subsection*{Common conventions: Source mathematical conventions}

\textbf{Well-defined objects.}
All probability spaces and measurable variables are specified or inherited
from a local statistical model. Expectations used as finite targets require
the stated integrability. A decision rule or estimator is measurable with
respect to its available information; randomized rules may use an auxiliary
independent random variable. Norms without a subscript are Euclidean in
finite-dimensional spaces. An optimizer is requested only when attainment
is assured; otherwise the question concerns the optimal value and
$\varepsilon$-optimal admissible rules for every $\varepsilon>0$.

\textbf{Identification.}
A structural model $m\in\mathcal M$ implies an observable law
$P_m^{\rm obs}$ and target $\theta(m)$. The target is point identified on
$\mathcal M$ if
\[
 P_{m_1}^{\rm obs}=P_{m_2}^{\rm obs}
 \quad\Longrightarrow\quad\theta(m_1)=\theta(m_2)
 \qquad(m_1,m_2\in\mathcal M).
\]
When this fails, the sharp identified set is
\[
 \Theta(P;\mathcal M)=\{\theta(m):m\in\mathcal M,
                                  P_m^{\rm obs}=P\}.
\]
Specifying an intervention or estimand does not establish identification.
Positivity, exclusion, causal sufficiency, measurement restrictions, and
transport assumptions are substantive local inputs, never automatic
consequences of notation.

\textbf{Minimax risk and nontrivial inference.}
For a statistical experiment with data law $P^{(n)}$, target $\theta(P)$,
and nonnegative loss $\ell$, define
\[
 \mathcal R_n(\mathcal P,\ell)
   =\inf_{\widehat\theta_n}\sup_{P\in\mathcal P}
      \E_{P^{(n)}}\ell(\widehat\theta_n,\theta(P)).
\]
The infimum is over measurable estimators. A sharp rate is a sequence
$r_n$ for which $0<c\leq C<\infty$, independent of $n$, satisfies
$c r_n\leq\mathcal R_n\leq C r_n$ for all sufficiently large $n$.
Squared-error parametric risk is of order $n^{-1}$; the corresponding
root-mean-squared error is of order $n^{-1/2}$. These different conventions
are distinguished locally. A uniform asymptotic confidence region $C_n$
must satisfy
\[
 \liminf_{n\to\infty}\inf_{P\in\mathcal P}
     \Pr_{P^{(n)}}\{\theta(P)\in C_n\}\geq 1-\alpha,
 \qquad 0<\alpha<1,
\]
together with an informative diameter or expected-width requirement.
Coverage alone permits the vacuous whole parameter space. Testing questions
similarly impose both size control and power against a declared separated
alternative class.

\textbf{Binary causal baseline.}
When an entry explicitly invokes this baseline, one observes independent
copies of $O=(X,W,Y)$, with $W\in\{0,1\}$ and potential outcomes $Y(0),Y(1)$.
Assume consistency $Y=Y(W)$, no interference, conditional exchangeability
$(Y(0),Y(1))\perp W\mid X$, and overlap
$\epsilon\leq e(X)\leq1-\epsilon$ almost surely for a fixed
$0<\epsilon<1/2$, where
\[
 e(x)=\Pr(W=1\mid X=x),\qquad
 m_w(x)=\E[Y\mid W=w,X=x],\qquad
 \tau(x)=m_1(x)-m_0(x).
\]
These assumptions identify conditional mean effects almost everywhere
under the covariate law. Pointwise or uniform effects require appropriate
regularity and a specified support. Continuous treatment, longitudinal
data, adaptive experiments, and network interference require their own
assumptions in place of this baseline. Bounded outcomes, densities, and
smoothness are additional assumptions only where stated.

\textbf{Function classes.}
On $[0,1]^d$, $\mathcal H_d^s(L)$ is the isotropic Hölder ball of radius
$L>0$: put $k=\lceil s\rceil-1$ and $\gamma=s-k\in(0,1]$, bound derivatives
through order $k$, and require the order-$k$ derivatives to be
$\gamma$-Hölder with constant at most $L$. Thus integer orders use a
Lipschitz final derivative convention. The class
$\operatorname{BV}_d(L)$ consists of functions $f\in L^1((0,1)^d)$ whose
distributional derivative is a finite vector measure and for which
$\|f\|_{L^1}+|Df|((0,1)^d)\leq L$.
An $L^\infty$ bound or a particular pointwise version is imposed separately
when needed. Sparse and neural-network classes require a declared
dictionary or architecture and coefficient bounds.

An anisotropic H\"older class with declared block smoothness indices means
coordinatewise H\"older regularity in each block, uniformly in the other
blocks, with all corresponding derivative and H\"older bounds at most
the stated radius. Derivative orders and the integer-order convention in
each block are those just given. Mixed-derivative bounds are additional
restrictions only when explicitly stated.

\textbf{Decision and computational questions.}
Policy regret compares admissible feasible policies under the same law and
constraints. In computational entries, finite data and thresholds have
rational encodings; running time is measured in their total bit length.
Real-valued equilibrium search uses the explicitly declared algebraic or
FIXP encoding. An approximation ratio for maximization is written
$\operatorname{OPT}/\operatorname{ALG}$ and for minimization as
$\operatorname{ALG}/\operatorname{OPT}$, with zero optima and infeasible
instances handled locally. Historical approximation bounds are not
automatically current bounds.
% END ECONOMICS STATEMENT
\end{document}
